\chapter{Discussione}
\label{cap:discussione}

Il capitolo confronta i risultati con i meccanismi esaminati e separa le
spiegazioni sostenute dalle diagnostiche dalle ipotesi ancora da verificare.

\section{Compressione dimensionale}
\label{sec:compressione}

In tutti e tre gli esperimenti che lavorano su dati ad alta dimensione il
circuito quantistico opera \emph{dopo} una compressione severa, e la compressione
è realizzata da una trasformazione classica:

\begin{center}
\begin{tabular}{@{}l l r@{}}
\toprule
\textbf{Esperimento} & \textbf{Compressione} & \textbf{Fattore} \\
\midrule
01 & $\demb = 8 \to \nq = 4$ per testa, \code{nn.Linear} addestrabile & $2\times$ \\
03 & $\dpatch = 147 \to \dvit = 10$, \code{nn.Linear} addestrabile & $14{,}7\times$ \\
04 & $784 \to 8$, PCA non addestrabile & $98\times$ \\
\bottomrule
\end{tabular}
\end{center}

La trasformazione successiva al collo di bottiglia può usare soltanto
l'informazione conservata dalla compressione. I confronti valutano quindi il VQC
nelle dimensionalità indicate in tabella, non prima della riduzione.

Nell'esperimento~01 la differenza media di validation loss è $+0{,}00543$, con IC
al 95\% $[-0{,}01040;+0{,}02126]$. H1 non è supportata; l'intervallo non equivale
a un test di equivalenza.

Nell'esperimento~04 la PCA è fissa e comune ai kernel. Le otto componenti
conservano in media il $76{,}54\%$ della varianza totale, una misura che non
specifica quanta parte della struttura discriminante rimanga nello spazio
compresso.

La dipendenza dalla compressione richiede confronti a dimensionalità diverse,
per esempio $\dvit \in \{4,10,16,32\}$ (\cref{sec:sviluppi}).

\section{Codifica angolare e saturazione}
\label{sec:saturazione}

La codifica angolare limita gli ingressi a $[-\pi,\pi]$ mediante una $\tanh$
(\cref{sec:quantum-linear}). Per $\abs{x}$ grande, la derivata tende a zero e
attenua il gradiente diretto allo strato di compressione.

Gli \emph{hook} dell'esperimento~03 registrano a ogni epoca saturazione e norma
media del gradiente (\cref{tab:diagnostica-03}). La saturazione media resta sotto
$0{,}013$ in ogni combinazione di dataset e variante. La norma del gradiente nel
modulo di compressione quantistico supera quella di \code{bounded\_mlp} nei tre
dataset. Le misure non indicano saturazione diffusa o gradienti più piccoli nel
modulo osservato, ma non coprono i parametri del circuito.

La connessione residua somma un ramo identità non limitato all'uscita del VQC,
confinata in $[-1,1]$. Un rapporto crescente fra le norme dei due rami
ridurrebbe il contributo relativo del circuito, ma gli \emph{hook} non registrano
questa quantità.

I \emph{barren plateau} riguardano la concentrazione dei gradienti nei circuiti
variazionali \parencite{mcclean2018barren}. Gli esperimenti usano 4--10 qubit e
due strati, senza misurare la varianza dei gradienti al variare di scala e
profondità. Le diagnostiche disponibili non verificano questo fenomeno.

\section{Capacità e sotto-adattamento}
\label{sec:capacita}

Il gap va letto insieme all'accuracy di training: può aumentare sia quando il
modello adatta meglio il training senza trasferire il guadagno al test, sia
quando cambiano entrambe le accuracy.

L'accuracy di training di \code{quantum\_reg} contro
\code{bounded\_mlp} è $0{,}950$ contro $0{,}886$ su PathMNIST, $0{,}916$ contro
$0{,}932$ su BloodMNIST e $0{,}832$ contro $0{,}770$ su DermaMNIST. Su PathMNIST e
DermaMNIST la variante quantistica adatta il training più strettamente della
baseline e ottiene un'accuracy di test inferiore; su BloodMNIST l'accuracy di
training è simile e quella di test resta inferiore. Il risultato non corrisponde
a un sotto-adattamento sistematico.

Il conteggio dei parametri non distingue una capacità effettiva ridotta da una
funzione difficile da ottimizzare. Curve ferme su un plateau basso favorirebbero
la prima spiegazione; curve ancora in salita o gradienti sistematicamente piccoli
la seconda. Le tracce raccolte non mostrano una delle due firme in modo stabile.

In nessuno dei tre dataset \code{quantum\_reg} riduce il gap rispetto a
\code{bounded\_mlp}; H3 non trova riscontro nella configurazione esaminata.

\section{Overhead di simulazione}
\label{sec:overhead}

Il rapporto fra i tempi della variante quantistica e di quella classica è
$46{,}5\times$ nell'esperimento 01, $9{,}29\times$ nel 02 e compreso fra
$60{,}7\times$ e $63{,}5\times$ nel 03; per il kernel dell'esperimento 04 il costo
è dominato dal calcolo quadratico delle similarità IQP. Il costo osservato dipende
dal backend a vettore di stato, dalla valutazione del QNode riga per riga
(\cref{sec:integrazione}), dal numero di qubit, dalla retropropagazione attraverso
le porte simulate e, nell'esperimento~01, dai
$n_{\text{layer}} \times \nhead \times 3$ circuiti valutati a ogni
\emph{forward pass}.

Il costo di simulazione ha imposto modelli piccoli, 4--10 qubit, cinque seed
negli esperimenti 02--03 e cento campioni nell'esperimento~04. I rapporti
temporali riguardano il simulatore e non l'esecuzione su hardware
(\cref{sec:costo}).

\section{Equità delle baseline}
\label{sec:equita-baseline}

Il conteggio dei parametri controlla solo una parte del confronto. La baseline
dell'esperimento~03 ha \num{1880} parametri di
\emph{patch embedding} contro \num{1540} della variante quantistica e ne replica
la limitatezza dell'output, mentre nell'esperimento~02 è la variante quantistica
ad averne 48 in più. Il numero di parametri non misura direttamente la capacità
delle due famiglie funzionali.

Learning rate, weight decay, numero di epoche e scheduler sono comuni e sono
stati fissati sulla baseline classica. Un tuning separato, con lo stesso budget
per ciascuna variante, rimane da eseguire (\cref{sec:sviluppi}).

Negli esperimenti 01 e 02 l'adapter classico addestrabile posto prima del
circuito rende l'effetto misurato quello della
composizione adapter${}+{}$VQC (\cref{sec:adapter}). L'esperimento~03 non usa
questo adapter. Nell'esperimento~04, IQP è confrontato con il coseno e con un RBF
a bandwidth mediana; non sono state provate altre famiglie di kernel o strategie
di selezione degli iperparametri.

\section{Limiti statistici}
\label{sec:limiti-statistici}

Il limite principale è la numerosità. Con cinque seed negli esperimenti 02 e 03 la
dimensione dell'effetto minima rilevabile con potenza $0{,}8$ è $\dz = 1{,}68$, e
con i dieci seed dell'esperimento 01 scende a circa $1{,}00$
(\cref{ssec:potenza}). Effetti moderati potevano quindi non essere rilevati. I
venti sottocampioni dell'esperimento~04 si sovrappongono e condividono il test
ufficiale nella validazione SVC; non sono venti dataset indipendenti.

Nell'esperimento~04 IQP aumenta $\rho$ rispetto al coseno ma riduce il criterio
di Fisher. $\rho$ dipende dalla scala e dal livello medio delle similarità; le
metriche della SVC forniscono il controllo predittivo. Balanced accuracy, AUC e
macro-F1 sono inferiori per IQP nei tre regimi di campionamento.

Non è stata definita una soglia di equivalenza né applicata una correzione per le
quattro ipotesi. L'analisi inferenziale riguarda le metriche primarie indicate
nel protocollo.

Nell'esperimento 01 la deviazione standard delle differenze di loss è
$0{,}02213$, contro una media di $0{,}00543$; nell'esperimento 02 la deviazione
delle differenze di accuracy è $0{,}0243$, contro una media in valore assoluto di
$0{,}0241$. Il segno della differenza varia quindi fra seed.

\section{Portata dei risultati negativi}
\label{sec:portata}

Le condizioni testate sono quelle riassunte nella \cref{sec:minacce}: codifica
angolare e ansatz \code{StronglyEntanglingLayers}, da 4 a 10 qubit e 2 strati
variazionali, in due sole posizioni architetturali, su MedMNIST a $28\times28$ e
su un corpus a livello di caratteri, con baseline appaiate e cinque o dieci seed.

In queste condizioni non emerge il miglioramento atteso della validation loss (H1) o
dell'accuracy di classificazione (H2); la macro-AUC di H2 è inferiore. Il gap di
generalizzazione non diminuisce ed è mediamente maggiore nei tre dataset di H3.
Per H4, il kernel di fedeltà IQP a 8 qubit ottiene valori inferiori al coseno e
all'RBF sia sul criterio di Fisher sia nella classificazione fuori campione.

Non sono stati esaminati altri ansatz, codifiche apprese, \emph{amplitude
embedding}, \emph{data re-uploading}, circuiti più grandi, altri domini o
hardware quantistico. Il lavoro non valuta vantaggi computazionali, cioè
l'affermazione (a) della \cref{ssec:quattro-promesse}.

I risultati riguardano le quattro configurazioni descritte. La variazione fra
seed rende instabile il segno delle differenze piccole.

L'infrastruttura comprende appaiamento dei seed, verifica dei componenti
condivisi, corruzioni deterministiche, manifest di provenienza e schema dei
risultati versionato. Può essere riutilizzata per altre configurazioni.
